% Generated draft: requires mathematical and source review.
\documentclass[11pt]{article}
\usepackage[margin=25mm]{geometry}
\usepackage{fontspec}
\setmainfont{Latin Modern Roman}
\usepackage{amsmath,amssymb,mathtools,unicode-math}
\setmathfont{Latin Modern Math}
\usepackage{enumitem,needspace,xurl,hyperref}
\setlength{\parindent}{0pt}
\setlength{\parskip}{5pt}
\setlength{\emergencystretch}{3em}
\newcommand{\N}{\mathbb{N}}
\newcommand{\Z}{\mathbb{Z}}
\newcommand{\Q}{\mathbb{Q}}
\newcommand{\R}{\mathbb{R}}
\newcommand{\C}{\mathbb{C}}
\newcommand{\F}{\mathbb{F}}
\newcommand{\A}{\mathbb{A}}
\newcommand{\PP}{\mathbb{P}}
\newcommand{\E}{\mathbb{E}}
\newcommand{\eps}{\varepsilon}
\newcommand{\dd}{\,\mathrm{d}}
\begin{document}
\section*{UnsolvedMath ID: 20002097}
\textbf{Source problem code:} catalogue entry
\subsection{Definitions and mathematical statement}
Polterovich's Question \(1.39\) conjectures that a contactomorphism \(\phi\) of a closed contact manifold whose induced contact-homology action is ``hyperbolic'' preserves no smooth volume form, calling this strongly dissipative. The source explicitly says that hyperbolic still needs a definition. A version of contact homology, coefficients, its automorphism and the intended hyperbolicity condition must therefore be supplied before this becomes a precise implication; no meaning is silently substituted from ordinary hyperbolic dynamics.

For the aid's cooriented formulation, let \(\dim M=2m+1\), choose \(\alpha\) with \(\xi=\ker\alpha\), and assume \(\phi\) preserves coorientation, so
\[
 \phi^*\alpha=e^g\alpha,\qquad
 \phi^*(\alpha\wedge(\dd\alpha)^m)
       =e^{(m+1)g}\alpha\wedge(\dd\alpha)^m.
\]
A smooth positive volume is \(e^f\alpha\wedge(\dd\alpha)^m\); it is invariant precisely when
\[
 (m+1)g=f-f\circ\phi.
\]
Equivalently \(g=h-h\circ\phi\) for some smooth \(h\), and this is equivalent to strictness for the conformally changed contact form \(e^h\alpha\). The intrinsic obstruction is the algebraic smooth dynamical-cohomology class
\[
 [g]_\phi\in C^\infty(M)/\{h\circ\phi-h:h\in C^\infty(M)\}.
\]
Changing \(\alpha\) to \(e^v\alpha\) changes \(g\) by \(v\circ\phi-v\), so this class is contact-form independent. The undefined contact-homology conjecture is thus reduced to proving that its eventual hyperbolicity hypothesis forces \([g]_\phi\ne0\).

A sufficient test is a periodic point \(x\), \(\phi^N x=x\), with
\[
 \sum_{j=0}^{N-1}g(\phi^j x)\ne0.
\]
Every coboundary has zero such sum. The converse is not asserted for a general contactomorphism without additional cohomological-dynamics hypotheses.

For the normalized cotangent lift of a diffeomorphism \(f:Q\to Q\), choose a Riemannian unit cosphere \(S^*Q\). The lift is
\[
 \Psi_f(q,p)=\left(f(q),\frac{(Df_q)^{-T}p}{\|(Df_q)^{-T}p\|}\right),
 \qquad \Psi_f^*\alpha=\frac{\alpha}{\|(Df_q)^{-T}p\|}.
\]
If \(q\) has period \(N\) and \((Df_q^N)^{-T}\) has a real eigen-covector \(p\) with eigenvalue \(\lambda\), \(|\lambda|\ne1\), then the corresponding unit covector is periodic for the normalized lift: after \(N\) if \(\lambda>0\), or after \(2N\) if \(\lambda<0\). Its conformal sum is respectively \(-\log|\lambda|\) or \(-2\log|\lambda|\), a nonzero certificate. This can happen even for a volume-preserving base diffeomorphism. The certificate is not claimed to prove that any specified contact-homology action is hyperbolic.

The source's broader theme is information retained by the smooth-to-contact constructions: a manifold \(Q\) gives its unit cotangent contact manifold, a suitable embedded submanifold \(K\) gives its unit conormal Legendrian, and a diffeomorphism gives a contact lift. It asks how much information is lost and refers to Eliashberg's first fifth-day question, absent here. Its final bracket asks how precisely this relates to the dissipation conjecture. That uncertainty is preserved; no general faithfulness or information-recovery theorem is supplied.
\subsection{Short English statement}
The source leaves contact-homology hyperbolicity undefined. The aid gives a precise target: absence of invariant volume is equivalent to a nonzero smooth conformal-factor cohomology class. Periodic conformal sums certify this, with an explicit real-eigen-covector test for cotangent lifts. The broader smooth-to-contact information question remains unspecified.
\subsection{Sources}
\begin{itemize}
\item Local UnsolvedMath registry, permanent ID \(20002097\), source code \texttt{catalogue entry}; canonical collection \texttt{aim-geometry-notes.json}, notes[434].
\item AIM workshop ``Holomorphic curves in contact geometry,'' Question \(1.39\), Polterovich's conjecture and theme: \url{https://aimath.org/WWN/contactgeom2/contactgeom2.pdf}. The supplied verification reports agreement of AIM HTML and PDF, with line-break repairs only; the hyperbolicity warning and bracketed questions are original. No external source or PDF was opened in this session.
\item The source refers to Eliashberg's first question on the fifth day without reproducing it. No full authoritative citations for the contact-homology setup, dynamical criterion or cotangent calculation are supplied. Original remarks and literature are empty, with no proposer-field value or year.
\item Supplied research attempt \(1\), result ``reduction,'' run status ``open,'' rigor label ``proved,'' novelty confidence ``low.'' Registry status \texttt{partially\_solved}, classification \texttt{PARTIAL-PROGRESS}, suggested difficulty \texttt{L4}, updated 14 August 2026.
\end{itemize}
\end{document}
